% a02_Kpos part deg2, continued: the tree near the face of the reduction cube (input from a02_deg2.tex after thm:a02:WR)
\paragraph{The tree near the face of the reduction cube.} Step (2) of the proof of \cref{thm:a02:WR} shows that the only source-free connected subgraph with $D\le0$ all of whose vertices are $\tilde A$-vertices is the tree $T$ of \cite[(2.4)]{Ba3}: two insertion vertices of \cite[(3.52)--(3.53)]{K1} of $J_1$ type joined by a doubly differentiated line, $D(T)=0$, no partner, inside family $\{T\}$. Its member-wise identity, a derivative at the end of a line (\cref{ssec:a02:ward}), is a sum over the whole lattice; at a fixed tuple the positions of insertion vertices are summed over the reduction cube $\square'$ with a sharp cut (\cite[p.~668]{K1}: only through the expansion \cite[(3.52)]{K1} can vertices approach $\partial\square'$), so near $\partial\square'$ the identity leaves a face term. The lines are decomposed scale by scale on $\square'$, $G_{(i)}(\square',0)=G_{(i)}(0)+\delta G_{(i)}(\square')$, the difference kept as an internal line of index $i$ (as in step (2) of the proof of \cref{thm:a02:WR}); we write $F_i(v)\defeq\exp[-\delta_1\dist(v,\partial\square')/\rho_i]$ with $\delta_1\defeq\delta_0/2$, $\delta_0$ the decay rate of the free slices \cite[(3.63)]{K1}.

\begin{lemma}[face factors]\label{lem:a02:FACE}
\textup{(a)} For $u,v\in\square'$ the line $\delta G_{(i)}(\square')(u,v)$ obeys the free slice bounds \cite[(3.63)--(3.65)]{K1} (sup, first differences, H\"older quotients of the first differences) at index $i$ with half the rate, times $F_i(u)F_i(v)$ (for a H\"older pair $u,\tilde u$, $F_i$ taken at the point of the pair nearer to $\partial\square'$). \textup{(b)} Let a term of the ordered bound of \cref{thm:a02:S3}(a) carry the factor $F_i(v)$ at a vertex $v$ of the tree $T$, $i$ the index of its line $\ell$. Then the term is an \textup{(R2)} unit in the sense of \cref{thm:a02:WR}: a credit $\rho_i$ on $\ell$ and a debit $\rho_{i'}^{-1}$ on one leg $l'$ of $T$, of index $i'\ge i$, or no debit when the surviving vertex of $T$ is the root of the graph; $D'(T)=1$, and $D'(\gamma)\ge D(\gamma)$ for every prefix component $\gamma\supsetneq T$.
\end{lemma}
\begin{proof}
(a) The boundary clause of \cite[Thm.~3.3]{K1} gives $\delta G_j(\square')$ and $\delta C^{(j)}(\square')$ their bounds with the extra factor $e^{-\delta_0\dist(\{u,v\},\partial\square')/\rho_j}$ (the clause used for the lines $\delta G_j(\square',A)$ in the proof of \cref{lem:a02:I} through \cite[(5.3)]{K2}). Writing each factor of the slice $G_{(i)}=a_iG_iQ^*C^{(i)}QG_i$ on $\square'$ as its free-lattice part plus its $\delta$-part and telescoping, every term of $\delta G_{(i)}(\square')$ carries one such factor, and the decay $e^{-\delta_0|\cdot|/\rho_i}$ of the factors (the $\delta$-factor's own included) moves it to both $u$ and $v$ at half the rate ($\min+\text{distance}\ge\max\ge\frac12(\dist(u,\partial\square')+\dist(v,\partial\square'))$), which is $F_i(u)F_i(v)$ with $\delta_1=\delta_0/2$ against the free bounds at the rate $\delta_0/2$; the H\"older clause follows from the sup and first-difference clauses of the factors as \cite[Prop.~3.7]{K1} follows from \cite[Thm.~3.3]{K1}, the boundary factor of the differentiated factor taken at the point of the pair nearer to $\partial\square'$, a pair at distance $>\rho_i$ being bounded by two first differences. On the $f_D$ flow the same follows from the image sum $\delta G_{(i)}(\square')(u,v)=\sum_{g\ne1}\operatorname{sgn}(g)\,G_{(i)}(0)(u,gv)$ of \cref{prop:a02:S}(c), since $|u-gv|\ge\max\{|u-v|,\dist(u,\partial\square')+\dist(v,\partial\square')\}$ for every $g\ne1$ and the images at distance $\le t$ number $O((1+t/p)^d)$. (b) In the ordered bound the line $\ell=(v,v')$ is shrunk by summing one endpoint against its kernel; $\dist(\cdot,\partial\square')$ is $1$-Lipschitz, so $F_i(v)\le e^{\delta_1|v-v'|/\rho_i}F_i(v')$ moves the factor to whichever endpoint survives, at the cost of half of the decay of $\ell$. The surviving vertex $x$ of $T$ is then summed against a leg $l'$ of $T$, a line of index $i'\ge i$ with the kernel profile $|K_{l'}(u-x)|\le C\rho_{i'}^{-c}e^{-\delta_0|u-x|/\rho_{i'}}$ of \cite[(3.63)--(3.65)]{K1} (a cut-line leg, \cref{lem:a02:XL}(b), has the unit-scale profile, $i'=j$), or it is the root. Since $F_i(x)\le\sum_\pi e^{-\delta_1\dist(x,\pi)/\rho_i}$ over the $2d$ face planes of $\square'$, each term depending on one coordinate only, summing that coordinate first gives $\sum_{x_\mu}5^{-j}e^{-\delta_1|x_\mu-\pi_\mu|/\rho_i}\le C\rho_i$ (as $5^{-j}\le\rho_i$), while the sum of $|K_{l'}|$ over the other $d-1$ coordinates is at most $C\rho_{i'}^{-c+d-1}$; so
\[
\sum_{x}5^{-3j}|K_{l'}(u-x)|\,F_i(x)\ \le\ C\,\frac{\rho_i}{\rho_{i'}}\,\rho_{i'}^{\,d-c},
\]
the bound $C\rho_{i'}^{d-c}$ of that step of the induction times $\rho_i/\rho_{i'}$. At the root, where the final position sum is paid by $|\square'|$ (\cite[p.~664]{K1}), $\sum_{x\in\square'}5^{-3j}F_i(x)\le C\rho_i\cdot(\text{area of }\partial\square')\le C(\rho_i/p)\,|\square'|$, the side of $\square'$ being $\ge4p$. In the telescoped product of the ordered bound (\cref{lem:a02:Dsp}(ii)) the factor $\rho_i/\rho_{i'}$ raises by one the exponent of exactly those prefixes that contain $\ell$ and not $l'$, so $D'(T)=D(T)+1=1$ and $D'(\gamma)\in\{D(\gamma),D(\gamma)+1\}$ for $\gamma\supsetneq T$; at the root the credit has no debit.
\end{proof}

\begin{proposition}[the tree at every tuple]\label{prop:a02:TREE}
Let $T$ be the tree \cite[(2.4)]{Ba3} as a prefix component at a fixed tuple, in either branch (with $\square'_K$ in place of $\square'$ in the branch $R>p$), with vertices $v_1$, the reference point $x$, and $v_2$, both summed over $\square'$, its line $\ell$ of index $i$, and the products $W(v_1)$, $W'(v_2)$ of the factors at its vertices (an $\tilde A$ field, the weight $\zeta$ or $1-\zeta$ of step (2) of the proof of \cref{thm:a02:WR}, and one chain-end line of index $\ge i$, a cut line in the branch $R>p$). Then $\sum_{v_1,v_2\in\square'}\ell(v_1,v_2)W(v_1)W'(v_2)$ is a sum of terms of the shapes \textup{(R0)}--\textup{(R2)} of \cref{thm:a02:WR}, with $D'(T)\ge\min\{\alpha,1\}>0$ and $D'(\gamma)\ge D(\gamma)$ for $\gamma\supsetneq T$. In particular the member-wise identity of the tree holds at every tuple, near $\partial\square'$ included, which is what \cref{cor:a02:DEGall}(a) uses for it.
\end{proposition}
\begin{proof}
If $\ell=D_{b_1}D_{b_2}G_{(i)}(0)$, write
\[
\sum_{v_1,v_2\in\square'}\ell\,W\,W'=\sum_{v_1\in\square'}W(v_1)\sum_{v_2\in\square'}\ell(v_1,v_2)\big[W'(v_2)-W'(v_1)\big]+\sum_{v_1\in\square'}W(v_1)W'(v_1)\Big[\sum_{v_2}\ell(v_1,v_2)-\sum_{v_2\notin\square'}\ell(v_1,v_2)\Big].
\]
The full sum over $v_2$ vanishes exactly, a lattice difference summed over the whole lattice (or the torus) being zero, the derivative at the end of a line: \textup{(R0)}. The first term is the transport step of \cite[(3.10)]{Ba3} with $n=0$: the factor $|v_1-v_2|^\alpha$ is absorbed by the decay of $\ell$ at the price $\rho_i^\alpha$, a credit on $\ell$, and the $\alpha$-H\"older quotient of $W'$ costs $\rho_{i'}^{-\alpha}$ on the chain-end line at $v_2$, which is H\"older at its own scale up to $\partial\square'$ (the free part by \cite[(3.65)]{K1}, the $\delta G$ part by \cref{lem:a02:FACE}(a), a cut line by \cref{lem:a02:XL}(b)), and $O(1)$ on the smooth factors $\tilde A$ (\cite[(3.45)]{K1}) and $\zeta$ ($|\nabla\zeta|\le4/p$): \textup{(R2)} with $D'(T)=\alpha$. In the cut term $v_1\in\square'$ and $v_2\notin\square'$, so $|v_1-v_2|\ge\dist(v_1,\partial\square')$ and half of the decay of $\ell$ gives $F_i(v_1)$; the sum over $v_2$ is then bounded by the kernel integral of $\ell$, as at every step of the ordered bound, and \cref{lem:a02:FACE}(b) makes the term an \textup{(R2)} unit with credit $1$. If $\ell=\delta G_{(i)}(\square')$, \cref{lem:a02:FACE}(a) gives $F_i(v_1)$ (the factor $F_i(v_2)$ is dropped) and \cref{lem:a02:FACE}(b) applies; no identity is used. Since $T$ has no partner and its inside family is $\{T\}$, the rewriting is an identity term by term, and the nesting of step (4) of the proof of \cref{thm:a02:WR} is unchanged. Away from $\partial\square'$ and on the torus this is Ba{\l}aban's treatment of \cite[(2.4)]{Ba3} in \cite[Prop.~2.1]{Ba3}; the cut term is the face term that the cut of the position sums produces.
\end{proof}
